\documentclass[11pt,a4paper]{article}
\usepackage[T1]{fontenc}
\usepackage[utf8]{inputenc}
\usepackage[english]{babel}
\usepackage{lmodern}
\usepackage{microtype}
\usepackage{geometry}
\geometry{margin=2.5cm}
\usepackage{hyperref}
\hypersetup{
    colorlinks=true,
    linkcolor=blue!70!black,
    urlcolor=blue!70!black,
    citecolor=blue!70!black
}
\usepackage{booktabs}
\usepackage{longtable}
\usepackage{enumitem}
\usepackage{xcolor}
\usepackage{titlesec}
\usepackage{fancyhdr}
\usepackage{graphicx}
\usepackage{listings}

\definecolor{codegray}{rgb}{0.95,0.95,0.95}
\definecolor{doctum}{rgb}{0.10,0.30,0.60}

\lstset{
  backgroundcolor=\color{codegray},
  basicstyle=\small\ttfamily,
  breaklines=true,
  frame=single,
  rulecolor=\color{gray!40},
  xleftmargin=1em,
  framexleftmargin=0.5em
}

\pagestyle{fancy}
\fancyhf{}
\fancyhead[L]{\textcolor{doctum}{\small\textbf{Doctum Consilium}}}
\fancyhead[R]{\small\nouppercase{\leftmark}}
\fancyfoot[C]{\thepage}

\titleformat{\section}{\large\bfseries\color{doctum}}{}{0em}{}[\titlerule]
\titleformat{\subsection}{\normalsize\bfseries\color{doctum!80}}{}{0em}{}


\begin{document}

\title{
  \textcolor{doctum}{\Large\textbf{MedVision AI}} \\[0.5em]
  \large Medical Image AI Pipeline — Brain MRI Classification and Segmentation \\[0.3em]
  \normalsize Learning Resources and Contributor Onboarding
}
\author{Doctum Consilium \\ \small\textit{Generated 2026-05-08}}
\date{}
\maketitle
\thispagestyle{fancy}

\begin{abstract}
This document provides a curated list of courses (Coursera, Pluralsight, A Cloud Guru),
training repositories, and reference books to help contributors ramp up on the
technologies used in \textbf{MedVision AI}.
Resources are selected for quality, relevance, and practical applicability.
\end{abstract}

\tableofcontents
\newpage

\section{Prerequisites}
Before starting the courses below, ensure you are comfortable with:
Python, linear algebra (matrices, eigenvalues), calculus (gradients), basic statistics.

\section{Coursera Courses}
\begin{longtable}{p{6cm} p{1.5cm} p{7cm}}
\toprule
\textbf{Course} & \textbf{Link} & \textbf{Why relevant} \\
\midrule
\endhead
  AI for Medicine Specialization & \href{https://www.coursera.org/specializations/ai-for-medicine}{\texttt{link}} & Medical image analysis, diagnostic AI, treatment effect estimation — directly applicable. \\
  Deep Learning Specialization & \href{https://www.coursera.org/specializations/deep-learning}{\texttt{link}} & CNN architecture theory, optimisation, regularisation. \\
  Practical Data Science on AWS & \href{https://www.coursera.org/specializations/practical-data-science}{\texttt{link}} & MLOps, experiment tracking, model deployment patterns. \\
\bottomrule
\caption{Coursera}
\end{longtable}


\section{Pluralsight Courses}
\begin{longtable}{p{6cm} p{1.5cm} p{7cm}}
\toprule
\textbf{Course} & \textbf{Link} & \textbf{Why relevant} \\
\midrule
\endhead
  PyTorch: Getting Started & \href{https://www.pluralsight.com/courses/pytorch-getting-started}{\texttt{link}} & Tensors, autograd, DataLoader — foundations for this project. \\
  MLflow for ML Experiment Tracking & \href{https://www.pluralsight.com/courses/mlflow-ml-experiment-tracking}{\texttt{link}} & Direct tool match. \\
\bottomrule
\caption{Pluralsight}
\end{longtable}


\section{A Cloud Guru Courses}
\begin{longtable}{p{6cm} p{1.5cm} p{7cm}}
\toprule
\textbf{Course} & \textbf{Link} & \textbf{Why relevant} \\
\midrule
\endhead
  Machine Learning on AWS & \href{https://acloudguru.com/course/machine-learning-on-aws}{\texttt{link}} & SageMaker patterns for scalable model training. \\
\bottomrule
\caption{A Cloud Guru}
\end{longtable}


\section{GitHub Training Repositories}
\begin{itemize}[leftmargin=1.5em]
  \item \href{https://github.com/Project-MONAI/MONAI}{\textbf{Project-MONAI/MONAI}} — Medical image deep learning framework; transforms, datasets, metrics.
  \item \href{https://github.com/mlflow/mlflow}{\textbf{mlflow/mlflow}} — Experiment tracking tool used in this project.
  \item \href{https://github.com/iterative/dvc}{\textbf{iterative/dvc}} — Data version control; pipeline stages, remote storage.
  \item \href{https://github.com/pytorch/vision}{\textbf{pytorch/vision}} — DenseNet, ResNet, Swin-V2 implementations.
\end{itemize}

\section{Reference Books}
\begin{itemize}[leftmargin=1.5em]
  \item \href{https://www.elsevier.com/books/deep-learning-for-medical-image-analysis/zhou/978-0-12-810408-8}{\textit{Deep Learning for Medical Image Analysis}} — Zhou et al.. Segmentation, classification, and detection in radiology.
  \item \href{https://www.oreilly.com/library/view/programming-pytorch-for/9781492045342/}{\textit{Programming PyTorch for Deep Learning}} — Ian Pointer. Practical PyTorch patterns for transfer learning.
\end{itemize}

\section{Suggested Learning Path}
\begin{enumerate}
  \item Complete prerequisites (1--2 weeks).
  \item Work through the Coursera specialisation most relevant to the repo's domain (4--8 weeks).
  \item Clone the repository, read \texttt{README.md} and \texttt{INFRASTRUCTURE.md}.
  \item Run the service locally following the \textbf{Local Run} section of \texttt{INFRASTRUCTURE.md}.
  \item Study the Pluralsight courses for the specific tools used (2--4 weeks, parallel).
  \item Contribute a small fix or documentation improvement as a first PR.
  \item Study the GitHub training repositories for advanced patterns.
\end{enumerate}

\section{Free Resources}
\begin{itemize}
  \item \href{https://docs.python.org/3/tutorial/}{\textbf{Python Official Tutorial}} — free, authoritative.
  \item \href{https://kubernetes.io/docs/tutorials/}{\textbf{Kubernetes Interactive Tutorials}} — browser-based k8s labs.
  \item \href{https://fastapi.tiangolo.com/tutorial/}{\textbf{FastAPI Tutorial}} — official, comprehensive.
  \item \href{https://pytorch.org/tutorials/}{\textbf{PyTorch Tutorials}} — official, covers all levels.
  \item \href{https://www.deeplearning.ai/short-courses/}{\textbf{DeepLearning.AI Short Courses}} — free, recent AI topics.
\end{itemize}

\end{document}
